% =====================================================================
%  ODI 2026 (NeurIPS 2026 Workshop on On-Device Intelligence:
%  Foundation Models under Real-World Constraints, Sydney) -- PAPERS TRACK
%  Limit: 5 pages main text, references and appendix excluded; the
%  appendix goes in the SAME pdf after the references.
%  Deadline: Aug 29, 2026 AoE  |  OpenReview:
%  https://openreview.net/group?id=NeurIPS.cc/2026/Workshop/ODI
%  DOUBLE-BLIND. No author names, no repository URLs, anonymised appendix.
%
%  BUILD:  pdflatex voicelink_odi.tex  (twice)   -- see README.md
%  Every number in this file is taken from docs/VERIFIED_NUMBERS.md
%  (2026-08-23) or mac/intel/BENCHMARKS.md; see README.md for the map.
% =====================================================================

\documentclass{article}

\usepackage{neurips_2026}          % anonymous by default; [final] for camera-ready

\usepackage[utf8]{inputenc}
\usepackage[T1]{fontenc}
\usepackage{hyperref}
\usepackage{url}
\usepackage{booktabs}
\usepackage{amsmath}
\usepackage{amssymb}
\usepackage{graphicx}
\usepackage{microtype}
\usepackage{xcolor}

% Drafting helper -- delete every use before submission.
\newcommand{\todo}[1]{\textcolor{red}{[\textbf{TODO:} #1]}}
\newcommand{\fn}[1]{\texttt{#1}}

\title{Fifteen Milliseconds to the Board: Local, Diversity-First Vocabulary
Selection for Wearable Augmentative Communication}

\author{Anonymous Author(s)\\Affiliation\\\texttt{email}}

\begin{document}
\maketitle

\begin{abstract}
Augmentative and alternative communication (AAC) boards give nonspeaking people
a grid of words to tap, but the words are programmed in advance and are
routinely wrong for the moment. We describe a wearable system that senses the
surrounding conversation and scene with local speech and vision--language
models and rebuilds a ten-word board in real time, entirely on a single local
host, with no data leaving the device. Three design decisions make it
deployable under real-world constraints. (1)~Every model call passes through
one fixed-signature function, so the inference backend is a configuration flag:
we swapped PyTorch for OpenVINO without touching firmware or UI. (2)~The board is
composed for \emph{coverage} rather than likelihood: maximal-marginal-relevance
selection over a generated pool raises distinct communicative functions per
board from 3.88 to 4.68, cuts near-duplicate pairs $69\times$ (0.0274
$\rightarrow$ 0.0004) and raises coverage-of-intent from 0.78 to 1.00 over 120
scenarios. (3)~We distil the expensive candidate generator into a 769-d MLP but
keep the selection policy symbolic; the raw student argmax board was measured
useless, and re-running the unchanged MMR over the student's logits recovered
94\% of the teacher's functional diversity. The student exports to OpenVINO IR
($1.4\times10^{-6}$ logit parity) and NNCF INT8 ($2.8\times$ smaller, top-6
overlap 0.950 mean / 0.667 min). A profile of the end-to-end path found that 71
of 73\,ms per board was re-embedding a static 94-word vocabulary; memoising it
gives a 15.3\,ms p50 board on the fully-OpenVINO path. We report a diversity gate
that still fails, an evaluation that is entirely synthetic and LLM-judged, and
the absence of any Intel-hardware measurement.
\end{abstract}

% =====================================================================
\section{Introduction}
% =====================================================================

An AAC board is a grid of words and pictograms that a person who cannot speak
taps to communicate. Its vocabulary is chosen in advance by a caregiver or
speech--language pathologist who is guessing what the user will need; when a
teacher asks ``what do plants take in from the air?'' no pre-programmed board
has \emph{carbon dioxide} on it, and changing the board takes another person and
time. The system studied here removes the guess: a body-worn capture device
streams audio and low-rate images to a nearby computer, which transcribes the
conversation, describes the scene, combines both with a private personal-context
store, and pushes a board that fits \emph{this moment} to a tablet.

The setting imposes the constraints this workshop is about. Inference must be
local: the device continuously captures bystanders, so nothing may leave the
host. It must be fast: a board that arrives after the conversational turn has
moved on is not a degraded answer but no answer. It must be small enough to
run beside speech and vision models on a consumer host, and it must be
\emph{portable} across hosts, because the development machine (Apple ARM) is
not the deployment target (an Intel AI PC via OpenVINO). And it must be
\emph{reliable} in a specific sense: the user gets one slate and cannot scroll
or retry, so a board of seven paraphrases of the likeliest answer is worthless
even when the likeliest answer is correct. Prior context-aware AAC work
optimises keystroke savings \citep{kwickchat2022,praact2024}, the right
objective when composing a sentence and the wrong one when catching a turn.

This paper is a systems report: the topology and seam (\S\ref{sec:system}),
diversity-first slate composition and its metrics (\S\ref{sec:slate}), hybrid
distillation and the ONNX $\rightarrow$ OpenVINO IR $\rightarrow$ NNCF INT8 path
(\S\ref{sec:distil}), end-to-end latency engineering (\S\ref{sec:latency}), and
limitations we consider load-bearing, including a gate that fails
(\S\ref{sec:limits}).

% =====================================================================
\section{System: three nodes, one seam}\label{sec:system}
% =====================================================================

\textbf{Topology.} A wearable capture device (ESP32-S3 with an OV2640-class
camera and an I$^2$S MEMS microphone) senses and streams: one VGA JPEG every
2\,s and continuous 16\,kHz mono PCM, type-tagged and multiplexed over a single
WebSocket. It performs \emph{no} inference and hosts its own access point, so it
needs no router. A host computer runs all perception and ranking. A tablet or
phone renders the board as a single-origin web page over plain HTTP with a
WebSocket for live updates, and reports every tap. The wearable talks only to
the host; the tablet talks only to the host; nothing leaves the local network.
On the device, the I$^2$S read-and-send loop is a high-priority FreeRTOS task
pinned to one core while camera capture and WebSocket servicing run on the
other, so a slow JPEG grab can never stutter audio. A physical button toggles
capture (the hardware kill switch) and an LED is lit whenever capture is active,
as the visible signal to bystanders.

\textbf{Perception.} Whisper \fn{base.en} \citep{whisper} (INT8) transcribes utterances
segmented by an energy VAD; a near-field loudness heuristic labels speech as
\emph{partner} versus \emph{ambient}, and detected questions raise an
\emph{addressed} flag. Qwen2.5-VL-3B \citep{qwen25vl} produces a one-sentence
scene description, but only when a 64$\times$64 grayscale change detector
(0.21\,ms per frame) reports that the scene moved; VLM latency is 0.93\,s at 256\,px
and 2.05\,s at 512\,px on the development host, so the gate is what keeps the
audio path from blocking on vision. All model calls run in worker threads.

\textbf{The seam.} Every model call in the system passes through one function
with a fixed signature,
\begin{center}
\fn{rank\_vocabulary(scene, transcript, context, strategy) -> list[RankedWord]}.
\end{center}
The firmware never calls a model; the web page never calls a model; the server
reaches intelligence only through this function. The \fn{strategy} argument
selects the policy (\fn{topk}, \fn{diverse}, \fn{student}); environment flags
select the runtime beneath it (\fn{torch} or \fn{openvino} for the student
head, for the embedder, and the OpenVINO device). We make no research claim for
this; it is an engineering constraint. But it is what allowed the OpenVINO port
in \S\ref{sec:distil} to be a configuration change rather than a rewrite, it is
what lets a single benchmark harness time three runtime combinations through
the identical code path (\S\ref{sec:latency}), and it is what makes the latency
figures attributable to one replaceable component. A vocabulary-hash sidecar on
each exported artifact makes a stale model refuse to serve the wrong words.

\textbf{Privacy by architecture.} Audio and frame buffers are ephemeral and
never written to disk; the VLM emits one sentence and detects no faces; the
personal-context store never leaves the host; the only persisted data is the
user's own tap log. Four safety words are pinned outside model control: the
model suggests, the tap is the utterance.

% =====================================================================
\section{Diversity-first slate composition}\label{sec:slate}
% =====================================================================

The obvious objective for a context-aware board---rank candidate words by
likelihood, show the top $k$---is wrong for an interface that gives the user
one slate with no recovery. We found early that the obvious metric (does the
board contain the oracle word?) is easy to game: a board of ten synonyms of
the likeliest answer scores well and is useless. Our production policy
therefore selects for coverage of the communicative moves a user might want to
make, and our harness measures slates rather than words.

\textbf{Policy.} A local LLM over-generates a candidate pool with a slot-based
prompt: exactly one first-person phrase for each of nine AAC communicative
functions (request, reject, comment, question, share-information,
express-feeling, social, direct-attention, clarify) plus phrases on the concrete
topics in play. A deterministic core pool is injected so a bad generation can
never strand the user. The displayed slate is selected by maximal marginal
relevance \citep{mmr1998} over MiniLM \citep{minilm} sentence embeddings,
\begin{equation}
s^{*} = \arg\max_{s \in C \setminus S}
\Big[\lambda\,\mathrm{rel}(s, c) - (1-\lambda)\max_{s' \in S}\mathrm{sim}(s, s')\Big],
\end{equation}
with per-candidate function tags so selection spans functions, a lower
$\lambda$ when a question is active so answers spread across distinct responses
rather than paraphrases, an answer-quota reservation, and a hedge cap so
``I'm not sure'' cannot fill the board. Four pinned core words plus six dynamic
words give a ten-word board; metrics judge the top seven.

\textbf{Metrics.} Each candidate is assigned a function by multi-exemplar
nearest-centroid classification in embedding space (a single-centroid version
mislabelled both \emph{yes} and \emph{carbon dioxide} as \emph{I don't know};
fixing it changed the measured outcome, and the fix was applied to both arms).
We report \emph{distinct functions@7}, the \emph{near-duplicate pair rate}
(cosine above a threshold), the \emph{same-function-pair rate} (the direct
measure of ``seven paraphrases''), and \emph{coverage-of-intent}: the fraction
of scenarios in which at least one authored \emph{acceptable move} is
expressible from the board. Definitions are in Appendix~\ref{app:metrics}.

% Figure dropped for the 5-page limit; re-enable if space allows.
% \begin{figure}[t]
% \centering
% \includegraphics[width=0.38\linewidth]{../../mac/evaluation/phase3/figures/phase3_distinct_functions.png}\hfill
% \includegraphics[width=0.38\linewidth]{../../mac/evaluation/phase3/figures/phase3_near_dup.png}
% \caption{Diversity-first (\fn{diverse}) versus likelihood-ranked (\fn{topk})
% selection over 120 scenarios: distinct communicative functions in the top seven
% (left, higher is better) and near-duplicate pair rate (right, lower is better).}
% \label{fig:phase3}
% \end{figure}

\begin{table}[t]
\centering
\small
\caption{Diversity-first versus likelihood-ranked selection, 120 synthetic
scenarios across five settings (classroom, clinic, airport, dinner, direct
question), judged at $k{=}7$.}
\label{tab:selection}
\begin{tabular}{lcc}
\toprule
 & \textbf{Diversity-first (MMR)} & \textbf{Top-$k$} \\
\midrule
Distinct communicative functions per board ($\uparrow$) & \textbf{4.68} & 3.88 \\
Near-duplicate pair rate ($\downarrow$)                 & \textbf{0.0004} & 0.0274 \\
Coverage-of-intent ($\uparrow$)\footnotemark            & \textbf{1.00} & 0.78 \\
Exact match to oracle word ($\uparrow$, diagnostic)      & \textbf{0.525} & 0.333 \\
\bottomrule
\end{tabular}
\end{table}
\footnotetext{Re-scored on 2026-08-23 with the current scorer. The committed
result file, written in June before a word-boundary fix to the scorer, reports
0.81 for top-$k$ (and exact match 0.567 vs.\ 0.342); the cached boards are
unchanged. We quote the current scorer throughout.}

\textbf{Result.} Table~\ref{tab:selection} compares
the policy with a top-$k$ baseline that ranks the same pool by relevance. Every
measure improves, including exact match to the oracle word, which the baseline
is the one optimising for: on this evaluation, spreading the slate across
functions is not a trade against accuracy. The near-duplicate rate falls
$69\times$ (0.02738/0.000397). MMR is not new; the claim is that its objective is
the right one for a one-slate, no-recovery interface, and that slate-level
metrics are what an interactive deployment benchmark should report.

% =====================================================================
\section{Hybrid distillation and the deployment path}\label{sec:distil}
% =====================================================================

The teacher pipeline of \S\ref{sec:slate} produces a good board in 6.71\,s (IQR
6.63--6.97, 8 samples), an order of magnitude past any turn window. We distil
it, but replace only the expensive part.

\textbf{Learn relevance, keep selection symbolic.} The student is a
769-dimensional MLP: a 384-d MiniLM scene embedding, a 384-d transcript
embedding and a one-bit \emph{is-question} flag, through
LayerNorm $\rightarrow$ Linear(769$\rightarrow$192) $\rightarrow$ ReLU
$\rightarrow$ Dropout(0.3) $\rightarrow$ Linear(192$\rightarrow$94), trained
multi-label with sigmoid+BCE to predict which of 94 dynamic vocabulary words the
teacher placed on the board. The vocabulary is built by frequency from 600
teacher boards (min frequency 3, max 200, embedding-deduplicated at $\tau=0.80$
with 22 canonical merges); 559 rows survive deduplication and a 15\% category
cap, split 391/84/84 with a fingerprint so a drifted corpus refuses to score
against a stale partition. The student's raw argmax board was measured useless:
2.1 hedge tokens per top-7 and 72.6\% of boards with two or more hedges. The fix
was architectural rather than more training: treat the student's logits as a
candidate pool of $3k$ and run the \emph{unchanged} production MMR, answer quota
and hedge cap over it. Hedges fell to 0.012 per board (teacher 0.321), and the
device now runs exactly the policy the harness measures. The learned component
supplies relevance; the symbolic component supplies the diversity the argmax
discarded.

\begin{table}[t]
\centering
\small
\caption{Distilled student versus teacher on the held-out split ($n{=}84$;
$n{=}71$ for coverage, which needs authored acceptable moves). One gate fails
and is reported as failing.}
\label{tab:student}
\begin{tabular}{lccl}
\toprule
 & \textbf{Student} & \textbf{Teacher} & \textbf{Bar / verdict} \\
\midrule
Distinct functions@7                  & 4.083  & 4.333  & $\ge$0.90 retention: pass (0.942) \\
Near-duplicate pair rate              & 0.0006 & 0.0    & $\le$0.08: pass \\
Same-function-pair rate ($\downarrow$) & 0.1865 & 0.1684 & $\le$ teacher: \textbf{fail} \\
Board variety (77/84 distinct)        & 0.917  & ---    & $\ge$0.5: pass \\
Coverage-of-intent, dynamic slots ($n{=}71$) & 0.761 & 0.845 & $\ge$0.90$\times$: pass (0.901$\times$) \\
Coverage-of-intent, full board ($n{=}71$)     & 0.845 & 0.944 & reported only \\
Dynamic word agreement with teacher   & 0.056  & ---    & reported only (\S\ref{sec:limits}) \\
\bottomrule
\end{tabular}
\end{table}

Table~\ref{tab:student} gives the held-out result. The student retains 94.2\% of
the teacher's functional diversity and 90.1\% of its dynamic coverage-of-intent,
and fails one gate: it stacks slightly more same-function words than the
teacher. That gate has been reported as failing through three review rounds and
we are not rounding it away.

\textbf{ONNX $\rightarrow$ OpenVINO IR $\rightarrow$ NNCF INT8.} The student
exports to ONNX (opset 17, standard layers only) and converts to OpenVINO IR \citep{openvino}
with a maximum absolute logit delta of $1.43\times10^{-6}$ against ONNX over
random inputs. NNCF post-training quantisation on 360 real calibration contexts
yields an INT8 IR of 235\,KB against 661\,KB for FP32 ($2.8\times$), with top-6
overlap against FP32 of 0.950 mean and 0.667 minimum (gate: mean $\ge$0.8 and
min $\ge$0.5) and a maximum logit delta of 1.05 (mean 0.23). The MiniLM embedder
is exported with tokenizer, encoder and pooling in one graph, so the Intel
student path is torch-free end to end; it matches the PyTorch embeddings at
cosine 1.000 and encodes a single text in 5.6\,ms versus 16.0\,ms. A runtime
pitfall worth recording: OpenVINO's ARM CPU plugin defaults to f16 inference,
which made the Whisper decoder degenerate and drifted student logits by
$10^{-2}$; we pin \fn{INFERENCE\_PRECISION\_HINT=f32} on CPU.

% =====================================================================
\section{End-to-end latency under the seam}\label{sec:latency}
% =====================================================================

Micro-benchmarks of the student head alone are uninformative---0.015\,ms for
torch, 0.025\,ms for OpenVINO FP32 and 0.031\,ms for INT8 at batch~1, with
\fn{benchmark\_app} reporting 48k and 37k~FPS for the two IRs---because the
head is not where the time goes. We therefore time the whole call,
\fn{rank\_vocabulary(strategy=student)}, on the 84 held-out scenarios: MiniLM
embedding of scene and transcript, the head, and symbolic MMR selection, for
three runtime combinations selected purely by environment flags, with the
hardware fingerprint recorded alongside the result.

\textbf{The profile that mattered.} The committed Phase-4 measurement of this
path was 47.081\,ms per board (IQR 45.06--49.29, 50 samples) against the
6.71\,s teacher, a $142.6\times$ speedup. Profiling the call showed that 71 of
73\,ms per board was the MiniLM encoder re-embedding the static 94-word
vocabulary on every call so that MMR could compare candidates. Embeddings are
now memoised per (backend, text) and scene and transcript go through one batched
encoder call. The benchmark evicts the context texts between rounds, so the
vocabulary stays warm but every board still encodes a never-seen scene and
transcript, which is what a live device does.

\begin{table}[t]
\centering
\small
\caption{End-to-end board latency through the seam after memoisation, 84
held-out scenarios, Apple M4 host, OpenVINO 2026.2 ARM CPU plugin. ``Board
Jaccard'' is set overlap of the produced board with the torch path.}
\label{tab:latency}
\begin{tabular}{llllccc}
\toprule
\textbf{Backend} & \textbf{Embedder} & \textbf{Head} & \textbf{Device} & \textbf{p50 (ms)} & \textbf{p90 (ms)} & \textbf{Jaccard vs.\ torch} \\
\midrule
torch   & torch    & PyTorch        & cpu & 16.8 & 56.6 & 1.000 \\
ov-head & torch    & OpenVINO INT8  & CPU & 15.7 & 19.9 & 0.914 \\
ov-full & OpenVINO & OpenVINO INT8  & CPU & \textbf{15.3} & 27.4 & 0.914 \\
\bottomrule
\end{tabular}
\end{table}

Table~\ref{tab:latency} is the result: 15.3\,ms p50 on the fully-OpenVINO path
and 16.8\,ms on torch, from roughly 70\,ms before. Two points need care.
First, the $142.6\times$ figure is the \emph{committed} one, computed from the
47.081\,ms board of the earlier harness and a teacher timing that requires LLM
generation; the speedup has \emph{not} been re-measured against the teacher
since memoisation, and we do not quote a new ratio here. Second, the board the
OpenVINO path produces is not identical to the torch board (Jaccard 0.914),
which is the expected consequence of INT8 logits reordering a candidate pool at
the margin; whether this matters is a question for the slate metrics, not for
parity.

\textbf{Speech.} whisper-base.en INT8 through \texttt{openvino\_genai} runs
at 1162 and 1157\,ms p50 against 681 and 795\,ms for \texttt{faster-whisper}
on 2.1 and 4.0\,s clips ($n{=}10$, identical text) on this ARM host; an earlier
$n{=}1$ figure favouring OpenVINO did not reproduce. With the board at 15\,ms,
perception (0.7--1.2\,s speech, 1--4\,s scene) is the latency budget.

% =====================================================================
\section{Limitations}\label{sec:limits}
% =====================================================================

\textbf{All evaluation is synthetic and LLM-judged against an LLM teacher.}
The 120 Phase-3 scenarios and the 600-row teacher corpus are generated or
authored; function labels are assigned in embedding space and judged by an LLM
labelled \emph{proxy}; the teacher is itself an LLM. A student that passes here
may be imitating the teacher imitating a larger model. \textbf{There are zero
human subjects.} Because the wearable captures bystanders who have not
consented, testing is restricted to the authors by self-imposed policy until
written proxy consent and pseudonym handling exist. The tap log that would
replace synthetic evaluation is being collected from the first UI version but
currently holds 121 team-only rows and is not used for any number here.

\textbf{Zero Intel-hardware measurements.} Every latency and throughput figure
in this paper, including the OpenVINO ones, was measured on Apple ARM silicon
through the OpenVINO ARM CPU plugin. The OpenVINO path is validated for
correctness (parity, INT8 overlap, board Jaccard), not for speed on the
deployment target. The benchmark scripts self-label the hardware so that a run
on an Intel host is one command, but that run has not happened.
\todo{replace this paragraph if Intel (Tiber / Core Ultra) numbers land before the deadline}

\textbf{One gate fails.} Same-function-pair rate 0.1865 versus teacher 0.1684
(Table~\ref{tab:student}).

\textbf{Teacher word agreement is 0.056, and a closed vocabulary is why.} The
student picks from 94 words, and that vocabulary can represent only 32.5\% of
the teacher's dynamic word choices, 49.1\% of which are hedge/filler; 12.2\% of
rows have no representable dynamic word at all. A vocabulary $8.5\times$ larger
(800 words) still represents only 56.6\%. This is why we gate coverage-of-intent
rather than word agreement, and why concrete answers are now read out of the
context by an open-vocabulary extractor rather than from a label space. On the
held-out split that extractor fires on only 2 of 84 rows, so it moves no gate.
On a purpose-built set of 60 where/who/which/colour/alternative questions whose
answer is present in the scene or transcript, it fires on 88\% and puts an
expected answer on the board for 0.817 of scenarios versus 0.050 for the closed
vocabulary alone (`how many' questions: 0/6, by design). That set was authored
by us and is not independent evidence.

\textbf{Latency figures supersede one another.} The 47\,ms and 15.3\,ms boards
come from the harness before and after memoisation; we state which is which and
quote no speedup for the new one.

% =====================================================================
\section{Conclusion}
% =====================================================================

A context-aware AAC board can be built, compressed and ported to a second
runtime without a human-subjects study, an Intel machine, or a faithful
student---if the system is organised so those gaps stay visible. The seam makes
the backend swap a flag; slate metrics make ``diverse'' measurable; symbolic
selection lets a 235-KB INT8 student inherit the teacher's policy; timing the
whole call found the $4\times$ that mattered. What remains is consented
interaction data from the people the board is for.

% =====================================================================
\section*{Use of large language models}
% =====================================================================

\emph{Development:} LLM coding agents drafted code, tests and notes under human
direction and review; the authors own every design decision and claim, and each
number here was re-derived from committed result files. \emph{System under
study:} a local LLM is the teacher (\S\ref{sec:slate}), a VLM describes scenes,
and an LLM is the proxy judge; the consequences for validity are in
\S\ref{sec:limits}. LLM tools also edited this prose.

% =====================================================================
%  REFERENCES -- VERIFY EVERY ENTRY BEFORE SUBMITTING.
%  Entries marked VERIFY were carried from project notes and have not been
%  checked against the published record.
% =====================================================================
\begin{thebibliography}{9}

\bibitem[Carbonell and Goldstein(1998)]{mmr1998}
J.~Carbonell and J.~Goldstein.
\newblock The use of MMR, diversity-based reranking for reordering documents and
producing summaries.
\newblock In \emph{SIGIR}, 1998.

\bibitem[Qwen Team(2025)]{qwen25vl}
Qwen Team.
\newblock Qwen2.5-VL technical report.
\newblock \emph{arXiv preprint}, 2025.
% VERIFY: exact arXiv number and author list.

\bibitem[Wang et al.(2020)]{minilm}
W.~Wang, F.~Wei, L.~Dong, H.~Bao, N.~Yang, and M.~Zhou.
\newblock MiniLM: Deep self-attention distillation for task-agnostic compression
of pre-trained transformers.
\newblock In \emph{NeurIPS}, 2020.

\bibitem[Radford et al.(2023)]{whisper}
A.~Radford, J.~W. Kim, T.~Xu, G.~Brockman, C.~McLeavey, and I.~Sutskever.
\newblock Robust speech recognition via large-scale weak supervision.
\newblock In \emph{ICML}, 2023.
% VERIFY: confirm venue/year.

\bibitem[OpenVINO(2026)]{openvino}
Intel Corporation.
\newblock OpenVINO toolkit and the Neural Network Compression Framework (NNCF).
\newblock Software, release 2026.2, 2026.
% VERIFY: preferred citation form for OpenVINO / NNCF.

\bibitem[Shen et al.(2022)]{kwickchat2022}
J. Shen, B. Yang, J. J. Dudley, and P. O. Kristensson.
\newblock KWickChat: A multi-turn dialogue system for AAC using context-aware
sentence generation.
\newblock In \emph{IUI}, 2022.
% VERIFY: real author list and exact title -- carried from project notes.

\bibitem[Pereira et al.(2024)]{praact2024}
J. Pereira, L. Bittencourt, and others.
\newblock PrAACT: Predictive AAC using contextual transformers.
\newblock 2024.
% VERIFY: real venue, authors, title -- carried from project notes.

\end{thebibliography}

% =====================================================================
\appendix
% =====================================================================
\section{Metric definitions}\label{app:metrics}

All slate metrics are computed on the top seven words of a ten-word board (four
pinned core words plus six dynamic words) unless stated.

\begin{description}
\item[Distinct functions@7] Number of the nine AAC communicative functions
(request, reject, comment, question, share-information, express-feeling, social,
direct-attention, clarify) represented among the top seven words. Each word is
labelled by nearest-exemplar classification over multiple authored exemplars per
function in the 384-d MiniLM embedding space.
\item[Near-duplicate pair rate] Fraction of word pairs in the top seven whose
embedding cosine similarity is at or above a fixed threshold.
\item[Same-function-pair rate] Fraction of word pairs in the top seven that share
a function label. Lower is better; the Phase-4 bar is ``no worse than teacher''.
\item[Coverage-of-intent] Fraction of scenarios for which at least one of the
scenario's authored \emph{acceptable moves} is expressible from the board.
``Dynamic'' restricts the board to its non-core words and the moves to
non-hedge moves, because the four pinned core words self-certify on the full
board. Scored with word-boundary matching and a hedge module.
\item[Board variety] Number of distinct boards over the held-out set divided by
the set size; a degeneracy guard.
\item[Dynamic word agreement] Word-level overlap between the student's and the
teacher's dynamic words on the same scenario. Reported, not gated.
\item[Diversity retention] Student distinct-functions@7 divided by the teacher's.
\item[INT8 top-$k$ overlap] For each calibration context, $|T_k^{\mathrm{INT8}}
\cap T_k^{\mathrm{FP32}}|/k$ with $k{=}6$; mean and minimum over 360 contexts.
\item[Board Jaccard] Set overlap of the board produced by a backend with that of
the torch backend on the same scenario, averaged over the 84 held-out rows.
\end{description}

\section{Reproduction}\label{app:repro}

All commands run from the host server directory with perception disabled
(\fn{VOICELINK\_PERCEPTION=0}) and need no LLM, VLM or Whisper call unless
noted. \todo{anonymise or package the code before submission; this appendix
currently names scripts, not a repository}
\begin{small}
\begin{verbatim}
python3 -m pytest -q                                   # 144 passed, 0 xfailed
python3 evaluation/phase3/selftest.py --from-boards    # Table 1 (re-scores cached boards)
python3 evaluation/phase4/evaluate.py --teacher-samples 0   # Table 2 (student side)
python3 evaluation/phase4/evaluate.py --teacher-samples 8   # teacher latency (LLM run)
python3 -c "from intel import convert as c; print(c.verify_student_ir())"  # IR parity
python3 intel/quantize.py                              # NNCF INT8 + overlap gate
python3 intel/bench.py                                 # head micro-latency, benchmark_app
python3 intel/bench_e2e.py --n 84                      # Table 3 (end-to-end board)
python3 intel/bench_whisper.py --n 10                  # Whisper comparison
\end{verbatim}
\end{small}
Phase-4 training configuration hash: \fn{8151a7e34afd}; split fingerprint:
\fn{ab51f7fba4c9c39e}; vocabulary: 94 entries with 22 canonical merges.
Runtime selection: \fn{VOICELINK\_STUDENT\_BACKEND=torch|openvino},
\fn{VOICELINK\_EMBED\_BACKEND=torch|openvino}, \fn{VOICELINK\_OV\_DEVICE=CPU|GPU|NPU}.

\section{Hardware and software fingerprint}\label{app:hw}

\begin{itemize}
\item Host for every measurement: Apple M4 (arm64), Darwin 25.5.0. No Intel
hardware was used for any number in this paper.
\item OpenVINO 2026.2.1 (build 21919), ARM CPU plugin, devices [CPU],
\fn{INFERENCE\_PRECISION\_HINT=f32}; PyTorch 2.12.0; Python 3.13.6.
\item Head micro-benchmarks: batch 1, input (1, 769) float32, 500 timed
iterations after 50 warm-up; \fn{benchmark\_app} latency hint, 5\,s each.
\item End-to-end benchmark: 84 held-out scenarios, context texts evicted
between rounds, vocabulary embeddings warm.
\item Student artifacts on disk: torch 657\,KB, FP32 IR 661\,KB, INT8 IR
235\,KB (\fn{.xml}+\fn{.bin}).
\item Teacher: Qwen2.5-3B-Instruct (4-bit) via a local runtime; scene model
Qwen2.5-VL-3B-Instruct (4-bit); embedder \fn{all-MiniLM-L6-v2}; speech
Whisper \fn{base.en} INT8 via \fn{faster-whisper} (CTranslate2) or
\fn{openvino\_genai.WhisperPipeline}.
\item Wearable: ESP32-S3 (dual-core Xtensa LX7, 240\,MHz, 8\,MB PSRAM),
OV2640-class camera, SPH0645 I$^2$S microphone; VGA JPEG at 2000\,ms cadence,
16\,kHz PCM in 512-sample chunks. \todo{a bill of materials with prices does not
exist; do not quote a device cost}
\end{itemize}

\end{document}
